\documentclass[11pt]{article}

% Plantilla común para material de estudio y ayudantías.
% Documento A4, una columna, fuente Computer Modern de 11 puntos
% y márgenes de 2.5 cm por lado, siguiendo el formato usado
% en los desarrollos de Física Matemática II.

\usepackage[utf8]{inputenc}
\usepackage[T1]{fontenc}
\usepackage[spanish]{babel}
\usepackage{amsmath,amssymb,mathtools,bm,mathrsfs,graphicx,array,enumitem,caption,float,microtype}
\usepackage[dvipsnames]{xcolor}
\usepackage[a4paper,margin=2.5cm]{geometry}
\usepackage{tikz}
\usetikzlibrary{arrows.meta,calc,patterns,babel,decorations.pathreplacing,decorations.pathmorphing,positioning}
\usepackage[
    colorlinks=true,
    linkcolor=red,
    citecolor=green,
    urlcolor=red
]{hyperref}
\spanishdecimal{.}

\setlength{\parindent}{0pt}
\setlength{\parskip}{0.45em}
\setlength{\abovedisplayskip}{6pt}
\setlength{\belowdisplayskip}{6pt}
\setlength{\abovedisplayshortskip}{4pt}
\setlength{\belowdisplayshortskip}{4pt}
\setlist[itemize]{topsep=4pt,itemsep=2pt,leftmargin=1.8em}
\setlist[enumerate]{topsep=4pt,itemsep=2pt,leftmargin=1.8em}
\captionsetup{font=small,labelfont=bf}
\allowdisplaybreaks

% Colores usados en las figuras.
\definecolor{headinggray}{RGB}{55,55,55}
\definecolor{rulegray}{RGB}{150,150,150}
\definecolor{bluevec}{RGB}{29,78,216}
\definecolor{orangevec}{RGB}{180,83,9}
\definecolor{redvec}{RGB}{185,28,28}
\definecolor{greenvec}{RGB}{21,128,61}
\definecolor{fieldblue}{RGB}{30,90,180}
\definecolor{currentorange}{RGB}{195,95,20}
\definecolor{inducedgreen}{RGB}{20,120,80}
\definecolor{wiregray}{RGB}{55,55,55}
\definecolor{waveblue}{RGB}{29,78,216}
\definecolor{waveorange}{RGB}{180,83,9}
\definecolor{wavegreen}{RGB}{21,128,61}
\definecolor{wavered}{RGB}{185,28,28}

% Encabezado homogéneo.
\newcommand{\encabezado}[3]{%
{\Large\bfseries #1\par}
\vspace{2pt}
{\normalsize #2\hfill #3\par}
\vspace{6pt}\hrule\vspace{8pt}}

\newcommand{\problema}[2]{%
\vspace{0.55em}
\noindent\textbf{#1.}\enspace{\small #2}\par
\vspace{0.3em}}

\newcommand{\inciso}[1]{\vspace{0.35em}\noindent\textbf{(#1)}\enspace}

% Comandos auxiliares.
\newcommand{\dd}{\,\mathrm{d}}
\newcommand{\ii}{\mathrm{i}}
\newcommand{\e}{\mathrm{e}}
\newcommand{\sen}{\operatorname{sen}}
\newcommand{\grad}{\nabla}
\newcommand{\laplaciano}{\nabla^2}
\newcommand{\R}{\mathbb{R}}

\newcommand{\soluciontitulo}{%
  \par\medskip
  \noindent\textbf{Solución.}\par
  \smallskip
}

\newcommand{\resultadofinal}[1]{%
  \par\medskip
  \noindent\textbf{Resultado final.} #1\par
}

\newcommand{\paso}[1]{%
  \par\medskip
  \noindent\textbf{#1}\par
  \smallskip
}

% Entornos sin recuadro: se usan solo para integrar observaciones breves en la redacción.
\newenvironment{comentario}{\par\medskip\noindent}{\par\medskip}
\newenvironment{alerta}{\par\medskip\noindent\textbf{Observación.}\enspace}{\par\medskip}

\newcommand{\Bentra}[2]{%
    \draw[fieldblue,line width=0.55pt] (#1,#2) circle (0.11);
    \draw[fieldblue,line width=0.55pt] ($(#1,#2)+(-0.065,-0.065)$)--($(#1,#2)+(0.065,0.065)$);
    \draw[fieldblue,line width=0.55pt] ($(#1,#2)+(-0.065,0.065)$)--($(#1,#2)+(0.065,-0.065)$);
}
\newcommand{\Bsale}[2]{%
    \draw[fieldblue,line width=0.55pt] (#1,#2) circle (0.11);
    \fill[fieldblue] (#1,#2) circle (0.03);
}

\tikzset{
    >=Latex,
    eje/.style={->,thick,headinggray},
    onda/.style={very thick,waveblue},
    ondados/.style={very thick,waveorange},
    ondaaux/.style={thick,wavegreen},
    vec/.style={->,very thick,wavered},
    vecV/.style={->,very thick,bluevec},
    vecB/.style={->,very thick,orangevec},
    vecF/.style={->,very thick,redvec},
    vecE/.style={->,very thick,greenvec},
    panel/.style={draw=rulegray!80,fill=black!2,rounded corners=2pt},
    cargaPos/.style={circle,draw=bluevec,fill=blue!8,thick,minimum size=8mm,inner sep=0pt},
    cargaNeg/.style={circle,draw=greenvec,fill=green!8,thick,minimum size=8mm,inner sep=0pt},
    guia/.style={densely dashed,rulegray},
    etiqueta/.style={font=\footnotesize},
    nota/.style={draw=rulegray,fill=white,rounded corners=1.5pt,align=left,font=\footnotesize,inner xsep=6pt,inner ysep=5pt,line width=0.35pt}
}

\newcommand{\Fsale}[2]{%
    \draw[redvec,very thick] (#1,#2) circle (0.16);
    \fill[redvec] (#1,#2) circle (0.04);
}
\newcommand{\Fentra}[2]{%
    \draw[redvec,very thick] (#1,#2) circle (0.16);
    \draw[redvec,very thick] ($(#1,#2)+(-0.095,-0.095)$)--($(#1,#2)+(0.095,0.095)$);
    \draw[redvec,very thick] ($(#1,#2)+(-0.095,0.095)$)--($(#1,#2)+(0.095,-0.095)$);
}

\begin{document}

\encabezado{Fuerza magnética sobre cargas y conductores}{Campos y Ondas 510245}{J.R., F.B.}
\shorthandoff{<>}


\section*{Problema 1.}
Dibujar la fuerza magn\'etica sobre una carga el\'ectrica que entra a una regi\'on con campo magn\'etico uniforme. Se deben considerar los distintos casos de orientaci\'on de $\vec v$, $\vec B$ y signo de la carga.

\par\medskip\noindent La regla pr\'actica es siempre la misma: primero se calcula la direcci\'on de $\vec v\times\vec B$ como si la carga fuera positiva. Luego se aplica el signo de $q$. Si $q>0$ se conserva ese sentido; si $q<0$ se invierte. Por tanto, el criterio operacional consiste en aplicar de forma sistemática la regla de la mano derecha.\par\medskip 

\soluciontitulo

\par\medskip\noindent\textbf{Datos.}\par\smallskip
Los datos relevantes no son num\'ericos, sino geom\'etricos:
\begin{equation*}
    \vec F_B=q\,\vec v\times\vec B,
    \qquad
    |\vec F_B|=|q|vB\sin\theta.
\end{equation*}
Aqu\'i $\theta$ es el \~angulo entre $\vec v$ y $\vec B$.

\par\medskip\noindent\textbf{Modelo físico.}\par\smallskip
La fuerza magn\'etica es perpendicular a $\vec v$ y a $\vec B$. Por eso, para dibujarla, basta identificar el plano definido por ambos vectores y luego usar el producto cruz. Si $\vec v$ es paralela a $\vec B$, no hay componente perpendicular y la fuerza se anula.

\par\medskip\noindent\textbf{Análisis caso por caso.}\par\smallskip
Se trabaja con una base cartesiana directa, donde
\begin{equation*}
    \hat x\times\hat y=\hat z,
    \qquad
    \hat y\times\hat z=\hat x,
    \qquad
    \hat z\times\hat x=\hat y.
\end{equation*}

\textbf{Caso (a): $q>0$, $\vec v=v\hat x$ y $\vec B=-B\hat z$.}

\begin{center}
\begin{tikzpicture}[x=1cm,y=1cm,font=\small]
    \draw[panel] (-0.25,-0.35) rectangle (14.85,6.25);
    \node[font=\bfseries] at (7.30,5.86) {Caso (a): $q>0$, $\vec v=v\hat x$, $\vec B=-B\hat z$};

    \begin{scope}[opacity=0.62]
    \foreach \x in {0.9,2.2,3.5,6.7,8.0,9.3}{
        \foreach \y in {0.85,1.70,3.55,4.40}{
            \Bentra{\x}{\y}
        }
    }
    \end{scope}

    \node[cargaPos,minimum size=11mm,font=\large] (q) at (5.35,2.55) {$+$};
    \draw[vecV,line width=1.15pt] (q)--++(3.05,0) node[right] {$\vec v$};
    \draw[vecF,line width=1.15pt] (q)--++(0,1.70) node[above] {$\vec F_B$};

    \draw[->,thick] (0.55,0.35)--++(1.05,0) node[right] {$\hat x$};
    \draw[->,thick] (0.55,0.35)--++(0,1.05) node[above] {$\hat y$};
    \Fsale{0.55}{0.35}
    \node[anchor=west] at (0.86,0.95) {$+\hat z$ sale};

    \node[nota,anchor=west,text width=3.70cm] at (10.15,3.72)
    {Se calcula $\vec v\times\vec B$. Como $q>0$, la fuerza conserva ese sentido.};
    \node[etiqueta,anchor=west] at (10.25,0.95) {$\otimes$: $\vec B$ entrando al plano.};
\end{tikzpicture}
\captionof{figure}{Caso (a): la fuerza queda dirigida hacia $+\hat y$.}
\end{center}

\par\medskip\noindent\textbf{Cálculo.}\par\smallskip
\begin{align*}
    \vec v\times\vec B
    &=(v\hat x)\times(-B\hat z)\\
    &=-vB(\hat x\times\hat z)\\
    &=-vB(-\hat y)\\
    &=vB\hat y.
\end{align*}
Como $q>0$,
\begin{equation*}
    \vec F_B=qvB\hat y.
\end{equation*}

\textbf{Caso (b): $q<0$, $\vec v=-v\hat x$ y $\vec B=B\hat y$.}

\begin{center}
\begin{tikzpicture}[x=1cm,y=1cm,font=\small]
    \draw[panel] (-0.25,-0.35) rectangle (14.85,6.20);
    \node[font=\bfseries] at (7.30,5.82) {Caso (b): $q<0$, $\vec v=-v\hat x$, $\vec B=B\hat y$};

    \begin{scope}[opacity=0.72]
    \foreach \x in {1.05,2.45,3.85,8.10}{
        \draw[vecB,line width=0.82pt] (\x,0.65)--++(0,0.92);
        \draw[vecB,line width=0.82pt] (\x,3.48)--++(0,0.92);
    }
    \foreach \x in {1.05,2.45,8.10}{
        \draw[vecB,line width=0.82pt] (\x,2.05)--++(0,0.92);
    }
    \end{scope}

    \node[cargaNeg,minimum size=11mm,font=\large] (q) at (6.50,2.55) {$-$};
    \draw[vecV,line width=1.15pt] (q)--++(-3.00,0) node[left] {$\vec v$};

    \Fsale{6.50}{4.15}
    \node[redvec,anchor=west] at (6.90,4.15) {$\vec F_B$ sale};
    \Fentra{9.10}{1.45}
    \node[etiqueta,anchor=west,text width=4.30cm] at (9.55,1.45) {$\vec v\times\vec B$ entrar\'ia; como $q<0$, la fuerza se invierte.};

    \draw[->,thick] (0.55,0.35)--++(1.05,0) node[right] {$\hat x$};
    \draw[->,thick] (0.55,0.35)--++(0,1.05) node[above] {$\hat y$};
    \Fsale{0.55}{0.35}
    \node[anchor=west] at (0.86,0.95) {$+\hat z$ sale};
\end{tikzpicture}
\captionof{figure}{Caso (b): la carga negativa invierte el sentido de $\vec v\times\vec B$.}
\end{center}

\par\medskip\noindent\textbf{Cálculo.}\par\smallskip
\begin{align*}
    \vec v\times\vec B
    &=(-v\hat x)\times(B\hat y)\\
    &=-vB(\hat x\times\hat y)\\
    &=-vB\hat z.
\end{align*}
Como $q=-|q|$,
\begin{align*}
    \vec F_B
    &=(-|q|)(-vB\hat z)\\
    &=|q|vB\hat z.
\end{align*}
La fuerza sale del plano.

\textbf{Caso (c): $\vec v$ paralela o antiparalela a $\vec B$.}

\begin{center}
\begin{tikzpicture}[x=1cm,y=1cm,font=\small]
    \draw[panel] (-0.25,-0.35) rectangle (14.85,6.05);
    \node[font=\bfseries] at (7.30,5.68) {Caso (c): $\vec v\parallel\vec B$ o antiparalela a $\vec B$};

    \begin{scope}[opacity=0.72]
    \foreach \y in {0.85,1.65,3.45,4.25}{
        \draw[vecB,line width=0.82pt] (0.90,\y)--++(1.65,0);
        \draw[vecB,line width=0.82pt] (3.40,\y)--++(1.65,0);
        \draw[vecB,line width=0.82pt] (7.10,\y)--++(1.65,0);
    }
    \draw[vecB,line width=0.82pt] (0.90,2.55)--++(1.65,0);
    \draw[vecB,line width=0.82pt] (8.65,2.55)--++(1.65,0);
    \end{scope}

    \node[cargaPos,minimum size=11mm,font=\large] (q) at (6.50,2.55) {$+$};
    \draw[vecV,line width=1.15pt] (q)--++(-2.95,0) node[left] {$\vec v$};
    \node[draw=redvec,fill=white,rounded corners=2pt,text=redvec,font=\large,inner sep=6pt] at (11.70,1.65) {$\vec F_B=\vec0$};

    \draw[->,thick] (0.55,0.35)--++(1.05,0) node[right] {$\hat x$};
    \draw[->,thick] (0.55,0.35)--++(0,1.05) node[above] {$\hat y$};
    \node[nota,anchor=west,text width=3.55cm] at (10.15,3.82)
    {No hay componente perpendicular entre $\vec v$ y $\vec B$.};
\end{tikzpicture}
\captionof{figure}{Caso (c): si $\theta=0^\circ$ o $180^\circ$, entonces $\sin\theta=0$.}
\end{center}

\par\medskip\noindent\textbf{Cálculo.}\par\smallskip
Por ejemplo, si $\vec v=-v\hat x$ y $\vec B=B\hat x$, entonces
\begin{align*}
    \vec v\times\vec B
    &=(-v\hat x)\times(B\hat x)\\
    &=-vB(\hat x\times\hat x)\\
    &=\vec0.
\end{align*}
Equivalentemente,
\begin{equation*}
    |\vec F_B|=|q|vB\sin(180^\circ)=0.
\end{equation*}

\newpage

\textbf{Caso (d): campo magn\'etico inclinado.}

\begin{center}
\begin{tikzpicture}[x=1cm,y=1cm,font=\small]
    \draw[panel] (-0.25,-0.35) rectangle (14.85,6.55);
    \node[font=\bfseries] at (7.30,6.17) {Caso (d): campo inclinado, $\vec B=B_{\parallel}\hat y+B_{\perp}\hat x$};

    \begin{scope}[opacity=0.60]
    \foreach \x in {0.85,2.15,3.45}{
        \draw[vecB,line width=0.78pt] (\x,0.65)--++(0.78,0.78);
        \draw[vecB,line width=0.78pt] (\x,3.85)--++(0.78,0.78);
    }
    \foreach \x in {8.10,9.20}{
        \draw[vecB,line width=0.78pt] (\x,0.65)--++(0.78,0.78);
        \draw[vecB,line width=0.78pt] (\x,3.25)--++(0.78,0.78);
    }
    \end{scope}

    \node[cargaPos,minimum size=11mm,font=\large] (q) at (5.25,2.58) {$+$};
    \draw[vecV,line width=1.15pt] (q)--++(0,2.10) node[above] {$\vec v$};

    \draw[vecB,line width=1.15pt] (q)--++(2.45,2.00) node[above right] {$\vec B$};
    \draw[densely dashed,orangevec,line width=0.95pt] (q)--++(2.45,0) node[midway,below] {$\vec B_{\perp}$};
    \draw[densely dashed,orangevec,line width=0.95pt] ($(q)+(2.45,0)$)--++(0,2.00) node[midway,right] {$\vec B_{\parallel}$};

    \Fentra{8.85}{2.12}
    \node[redvec,anchor=west] at (9.20,2.12) {$\vec F_B$ entra};

    \draw[->,thick] (0.55,0.35)--++(1.05,0) node[right] {$\hat x$};
    \draw[->,thick] (0.55,0.35)--++(0,1.05) node[above] {$\hat y$};
    \Fsale{0.55}{0.35}
    \node[anchor=west] at (0.86,0.95) {$+\hat z$ sale};

    \node[nota,anchor=west,text width=3.75cm] at (10.25,4.55)
    {Solo $\vec B_{\perp}$ participa en el producto cruz.};
\end{tikzpicture}
\captionof{figure}{Caso (d): la componente perpendicular del campo fija la fuerza.}
\end{center}

\par\medskip\noindent\textbf{Cálculo.}\par\smallskip
Si $\vec v=v\hat y$ y $\vec B=B_{\parallel}\hat y+B_{\perp}\hat x$, entonces
\begin{align*}
    \vec v\times\vec B
    &=v\hat y\times(B_{\parallel}\hat y+B_{\perp}\hat x)\\
    &=vB_{\parallel}(\hat y\times\hat y)+vB_{\perp}(\hat y\times\hat x)\\
    &=\vec0-vB_{\perp}\hat z.
\end{align*}
Para $q>0$,
\begin{equation*}
    \vec F_B=-qvB_{\perp}\hat z.
\end{equation*}

\par\medskip\noindent\textbf{Resultado e interpretación física.}\par\smallskip
\begin{equation*}
    \boxed{(a)\;\vec F_B\text{ hacia }+\hat y},
    \qquad
    \boxed{(b)\;\vec F_B\text{ sale del plano}},
\end{equation*}
\begin{equation*}
    \boxed{(c)\;\vec F_B=\vec0},
    \qquad
    \boxed{(d)\;\vec F_B\text{ entra al plano}}.
\end{equation*}
La fuerza magn\'etica no apunta necesariamente en la direcci\'on del campo ni de la velocidad; apunta en la direcci\'on perpendicular determinada por el producto cruz.

\newpage

\section*{Problema 2.}
Una part\'icula cargada entra en un campo magn\'etico uniforme de m\'odulo $0.34\,\mathrm{mT}$ con rapidez $200\,\mathrm{m/s}$, perpendicular al campo. La carga tiene magnitud $3.2\times10^{-19}\,\mathrm C$. Calcular la fuerza magn\'etica que act\'ua sobre ella.

\par\medskip\noindent Como $\vec v\perp\vec B$, el factor angular es $\sin90^\circ=1$. Entonces toda la rapidez participa en la fuerza magn\'etica.\par\medskip 

\begin{figure}[h!]
\centering
\begin{tikzpicture}[x=1.18cm,y=1.18cm,font=\small]
    \draw[panel] (-3.85,-2.40) rectangle (4.25,2.55);
    \begin{scope}[opacity=0.58]
    \foreach \x in {-2.8,-1.45,-0.10,1.65,3.00}{
        \foreach \y in {-1.45,1.45}{\Bentra{\x}{\y}}
    }
    \foreach \x in {-3.15,3.35}{\Bentra{\x}{0.0}}
    \end{scope}
    \draw[bluevec,very thick] (0,0) circle (1.38);
    \node[cargaPos] (q) at (1.38,0) {$+$};
    \draw[vecV,line width=1.10pt] (q)--++(0,1.10) node[above] {$\vec v$};
    \draw[vecF,line width=1.10pt] (q)--++(-1.12,0) node[left] {$\vec F_B$};
    \node[nota,anchor=west,text width=2.42cm] at (1.80,1.10) {$\theta=90^\circ$\\$|\vec F_B|=|q|vB$};
    \node[etiqueta,anchor=west] at (-3.35,-1.95) {$\otimes$: $\vec B$ entra al plano};
\end{tikzpicture}
\caption{Caso perpendicular: la fuerza magn\'etica cambia la direcci\'on de la part\'icula.}
\end{figure}

\soluciontitulo

\par\medskip\noindent\textbf{Datos.}\par\smallskip
\begin{equation*}
    B=0.34\,\mathrm{mT},
    \qquad
    v=200\,\mathrm{m/s},
    \qquad
    |q|=3.2\times10^{-19}\,\mathrm C,
    \qquad
    \theta=90^\circ.
\end{equation*}
Convertimos el campo a unidades SI:
\begin{equation*}
    B=0.34\times10^{-3}\,\mathrm T=3.4\times10^{-4}\,\mathrm T.
\end{equation*}

\par\medskip\noindent\textbf{Modelo físico.}\par\smallskip
La magnitud de la fuerza magn\'etica sobre una carga en movimiento es
\begin{equation*}
    |\vec F_B|=|q|vB\sin\theta.
\end{equation*}

\par\medskip\noindent\textbf{Análisis.}\par\smallskip
Como la velocidad es perpendicular al campo, $\sin90^\circ=1$. Por lo tanto, la expresi\'on que se debe evaluar es
\begin{equation*}
    |\vec F_B|=|q|vB.
\end{equation*}

\par\medskip\noindent\textbf{Cálculos.}\par\smallskip
Sustituimos los datos directamente:
\begin{align*}
    |\vec F_B|
    &=(3.2\times10^{-19})(200)(3.4\times10^{-4})\,\mathrm N\\
    &=(3.2)(200)(3.4)\times10^{-23}\,\mathrm N\\
    &=2176\times10^{-23}\,\mathrm N\\
    &=2.176\times10^{-20}\,\mathrm N.
\end{align*}

\par\medskip\noindent\textbf{Resultado e interpretación física.}\par\smallskip
\begin{equation*}
    \boxed{|\vec F_B|\approx2.18\times10^{-20}\,\mathrm N}.
\end{equation*}
La fuerza es perpendicular a la velocidad; por eso modifica la direcci\'on del movimiento, no la rapidez instant\'anea de la part\'icula.

\newpage

\section*{Problema 4.}
En una regi\'on coexisten un campo el\'ectrico de $500\,\mathrm{N/C}$ y un campo magn\'etico de $0.15\,\mathrm T$, perpendiculares entre s\'i. Una part\'icula de carga $1\,\mathrm{mC}$ se mueve perpendicular a ambos campos y atraviesa la regi\'on sin desviarse. Calcular su velocidad.

\par\medskip\noindent Que la part\'icula no se desv\'ie no significa que no act\'uen fuerzas. Significa que la fuerza el\'ectrica y la fuerza magn\'etica se cancelan exactamente.\par\medskip 

\begin{figure}[h!]
\centering
\begin{tikzpicture}[x=1.16cm,y=1.16cm,font=\small]
    \draw[panel] (-3.90,-2.30) rectangle (4.05,2.95);
    \draw[greenvec,thick] (-3.35,1.48)--(3.35,1.48);
    \draw[greenvec,thick] (-3.35,-1.48)--(3.35,-1.48);
    \node[greenvec] at (-3.60,1.48) {$+$};
    \node[greenvec] at (-3.60,-1.48) {$-$};

    \foreach \x in {-2.65,-1.25,0.45,1.85,3.05}{
        \draw[vecE,line width=0.82pt] (\x,-1.10)--(\x,-0.34);
        \draw[vecE,line width=0.82pt] (\x,0.34)--(\x,1.10);
    }
    \begin{scope}[opacity=0.58]
    \foreach \x in {-2.65,-1.25,0.45,1.85,3.05}{
        \Bentra{\x}{-0.74} \Bentra{\x}{0.74}
    }
    \end{scope}

    \node[cargaPos] (q) at (-0.85,0) {$+$};
    \draw[vecV,line width=1.05pt] (-3.05,0)--(-1.20,0);
    \draw[vecV,line width=1.05pt] (-0.50,0)--(3.10,0) node[right] {$\vec v$};
    \draw[vecE,line width=1.05pt] (q)--++(0,1.02) node[above left] {$\vec F_E$};
    \draw[vecB,line width=1.05pt] (q)--++(0,-1.02) node[below left] {$\vec F_B$};
    \node[nota,anchor=center] at (1.75,2.38) {$\vec F_E+\vec F_B=\vec0$};
\end{tikzpicture}
\caption{Selector de velocidades: las fuerzas transversales se equilibran.}
\end{figure}

\soluciontitulo

\par\medskip\noindent\textbf{Datos.}\par\smallskip
\begin{equation*}
    E=500\,\mathrm{N/C},
    \qquad
    B=0.15\,\mathrm T,
    \qquad
    q=1\,\mathrm{mC}=10^{-3}\,\mathrm C.
\end{equation*}

\par\medskip\noindent\textbf{Modelo físico.}\par\smallskip
Act\'uan dos fuerzas:
\begin{equation*}
    |\vec F_E|=|q|E,
    \qquad
    |\vec F_B|=|q|vB,
\end{equation*}
porque $\vec v\perp\vec B$.

\par\medskip\noindent\textbf{Análisis.}\par\smallskip
La condici\'on de no desviaci\'on exige equilibrio transversal:
\begin{equation*}
    \vec F_E+\vec F_B=\vec0.
\end{equation*}
En magnitudes,
\begin{equation*}
    |\vec F_E|=|\vec F_B|.
\end{equation*}
Sustituimos las expresiones de cada fuerza:
\begin{equation*}
    |q|E=|q|vB.
\end{equation*}
Como $q\neq0$, la carga se cancela. Esto es clave: la velocidad seleccionada no depende del valor de la carga.

\par\medskip\noindent\textbf{Cálculos.}\par\smallskip
\begin{align*}
    |q|E&=|q|vB,\\
    E&=vB,\\
    v&=\frac{E}{B}.
\end{align*}
Sustituyendo los datos,
\begin{align*}
    v&=\frac{500}{0.15}\,\mathrm{m/s}\\
     &=3333.3\,\mathrm{m/s}.
\end{align*}

\par\medskip\noindent\textbf{Resultado e interpretación física.}\par\smallskip
\begin{equation*}
    \boxed{v\approx3.33\times10^3\,\mathrm{m/s}}.
\end{equation*}
El sistema deja pasar sin desviaci\'on solo a las part\'iculas con esta rapidez. Por eso esta configuraci\'on se interpreta como un selector de velocidades.

\newpage

\section*{Problema 6.}
Calcular la fuerza magn\'etica sobre una carga $q$ que entra en un campo magn\'etico de $4\,\mathrm T$ con rapidez $1500\,\mathrm{km/s}$, formando un \~angulo de $45^\circ$ con el campo magn\'etico.

\par\medskip\noindent\textbf{Observación.}\enspace El enunciado no entrega el valor num\'erico de $q$. Por eso el resultado debe quedar expresado en funci\'on de $|q|$.\par\medskip 

\begin{figure}[h!]
\centering
\begin{tikzpicture}[x=1.25cm,y=1.25cm,font=\small]
    \coordinate (O) at (0,0);
    \draw[panel] (-0.85,-1.35) rectangle (5.60,3.60);
    \draw[vecB,line width=1.05pt] (-0.45,0)--(4.85,0) node[right] {$\vec B$};
    \draw[vecV,line width=1.05pt] (O)--(2.85,2.85) node[above right] {$\vec v$};
    \draw[dashed,rulegray] (2.85,2.85)--(2.85,0);
    \draw[->,thick,greenvec] (O)--(2.65,0) node[midway,below] {$\vec v_{\parallel}$};
    \draw[->,thick,bluevec] (3.05,0)--(3.05,2.75) node[midway,right] {$\vec v_{\perp}$};
    \draw[thick] (0.72,0) arc (0:45:0.72);
    \node[etiqueta] at (1.03,0.38) {$45^\circ$};
    \Fsale{0.00}{-0.78}
    \node[redvec,etiqueta,anchor=west] at (0.38,-0.78) {$\vec F_B$ sale};
    \node[nota,anchor=west,text width=2.68cm] at (3.45,2.05) {$v_\perp=v\sin45^\circ$\\$|\vec F_B|=|q|Bv_\perp$};
\end{tikzpicture}
\caption{Solo la componente de la velocidad perpendicular a $\vec B$ produce fuerza magn\'etica.}
\end{figure}

\soluciontitulo

\par\medskip\noindent\textbf{Datos.}\par\smallskip
\begin{equation*}
    B=4\,\mathrm T,
    \qquad
    v=1500\,\mathrm{km/s}=1.5\times10^6\,\mathrm{m/s},
    \qquad
    \theta=45^\circ.
\end{equation*}
La carga queda como par\'ametro: $|q|$.

\par\medskip\noindent\textbf{Modelo físico.}\par\smallskip
La magnitud de la fuerza magn\'etica es
\begin{equation*}
    |\vec F_B|=|q|vB\sin\theta.
\end{equation*}
El factor $\sin\theta$ selecciona la parte perpendicular de la velocidad.

\par\medskip\noindent\textbf{Análisis.}\par\smallskip
Como $\theta=45^\circ$,
\begin{equation*}
    v_\perp=v\sin45^\circ=v\frac{\sqrt2}{2}.
\end{equation*}
Luego se puede usar $|\vec F_B|=|q|Bv_\perp$.

\par\medskip\noindent\textbf{Cálculos.}\par\smallskip
Primero calculamos la componente perpendicular:
\begin{align*}
    v_\perp
    &=(1.5\times10^6)\frac{\sqrt2}{2}\,\mathrm{m/s}\\
    &\approx1.06\times10^6\,\mathrm{m/s}.
\end{align*}
Entonces,
\begin{align*}
    |\vec F_B|
    &=|q|Bv_\perp\\
    &=|q|(4)(1.06\times10^6)\,\mathrm N\\
    &\approx4.24\times10^6|q|\,\mathrm N.
\end{align*}

\par\medskip\noindent\textbf{Resultado e interpretación física.}\par\smallskip
\begin{equation*}
    \boxed{|\vec F_B|\approx4.24\times10^6|q|\,\mathrm N}.
\end{equation*}
La componente paralela de la velocidad no contribuye al producto cruz. F\'isicamente, el campo solo curva el movimiento asociado a la componente perpendicular.

\newpage

\section*{Problema 9.}
Un conductor suspendido por dos cables tiene masa por unidad de longitud $\lambda=0.040\,\mathrm{kg/m}$. El campo magn\'etico es $B=3.6\,\mathrm T$ y entra en el plano de la hoja. Determinar la corriente necesaria para que la tensi\'on de los cables sea cero.

\par\medskip\noindent La condici\'on $T=0$ significa que los cables ya no sostienen el peso del conductor. Quedan como gu\'ia geom\'etrica, pero la fuerza vertical que equilibra al peso debe ser completamente magn\'etica.\par\medskip 

\begin{figure}[h!]
\centering
\begin{tikzpicture}[x=1.05cm,y=1.05cm,font=\small]
    \draw[panel] (-0.75,-1.55) rectangle (9.80,4.20);

    \begin{scope}[opacity=0.55]
    \foreach \x in {0.25,1.20,2.15,5.55,6.50,7.45}{
        \foreach \y in {-0.85,0.35,1.65,2.85}{\Bentra{\x}{\y}}
    }
    \end{scope}

    % Soporte superior fijo
    \draw[thick,black!70] (0.55,3.55)--(5.95,3.55);
    \foreach \x in {0.75,1.10,...,5.75}{
        \draw[black!45,line width=0.35pt] (\x,3.55)--++(-0.22,0.28);
    }
    \node[etiqueta,anchor=west] at (6.15,3.55) {soporte fijo};

    % cables de suspension
    \draw[gray!75,thick] (1.35,3.55)--(1.35,0.72);
    \draw[gray!75,thick] (5.15,3.55)--(5.15,0.72);
    \node[etiqueta,gray!70,anchor=east] at (1.25,2.35) {$T=0$};
    \node[etiqueta,gray!70,anchor=west] at (5.25,2.35) {$T=0$};

    % varilla/conductor
    \draw[line width=5.0pt,black!75,rounded corners=1.5pt] (1.05,0.72)--(5.45,0.72);
    \draw[vecV,line width=1.10pt] (1.50,1.08)--(3.15,1.08) node[midway,above] {$I$};

    \draw[vecF,line width=1.10pt] (3.25,0.95)--(3.25,2.40) node[left] {$\vec F_B$};
    \draw[->,very thick,black!70] (3.25,0.48)--(3.25,-1.00) node[left] {$\lambda L\vec g$};

    \node[nota,anchor=west,text width=2.65cm] at (6.55,1.35) {$\vec B$ entra al plano.\\Para que $T=0$, debe cumplirse $F_B=P$.};
\end{tikzpicture}
\caption{Conductor suspendido: los cables est\'an presentes, pero la condici\'on buscada es $T=0$.}
\end{figure}

\soluciontitulo

\par\medskip\noindent\textbf{Datos.}\par\smallskip
\begin{equation*}
    \lambda=0.040\,\mathrm{kg/m},
    \qquad
    B=3.6\,\mathrm T,
    \qquad
    g=9.8\,\mathrm{m/s^2}.
\end{equation*}
La corriente $I$ es la inc\'ognita.

\par\medskip\noindent\textbf{Modelo físico.}\par\smallskip
La fuerza magn\'etica sobre un conductor recto es
\begin{equation*}
    \vec F_B=I\vec L\times\vec B.
\end{equation*}
Como el conductor es perpendicular al campo,
\begin{equation*}
    |\vec F_B|=ILB.
\end{equation*}
El peso del conductor de largo $L$ es
\begin{equation*}
    P=(\lambda L)g.
\end{equation*}

\par\medskip\noindent\textbf{Análisis.}\par\smallskip
Si $T=0$, la fuerza magn\'etica debe sostener todo el peso:
\begin{equation*}
    |\vec F_B|=P.
\end{equation*}
Como ambas cantidades son proporcionales a $L$, conviene trabajar por unidad de longitud:
\begin{equation*}
    \frac{|\vec F_B|}{L}=IB,
    \qquad
    \frac{P}{L}=\lambda g.
\end{equation*}

\par\medskip\noindent\textbf{Cálculos.}\par\smallskip
Igualamos fuerza magn\'etica y peso:
\begin{align*}
    ILB&=\lambda Lg,\\
    IB&=\lambda g,\\
    I&=\frac{\lambda g}{B}.
\end{align*}
Sustituyendo los datos,
\begin{align*}
    I&=\frac{(0.040)(9.8)}{3.6}\,\mathrm A\\
     &=\frac{0.392}{3.6}\,\mathrm A\\
     &\approx0.1089\,\mathrm A.
\end{align*}

\par\medskip\noindent\textbf{Resultado e interpretación física.}\par\smallskip
\begin{equation*}
    \boxed{I\approx0.11\,\mathrm A}.
\end{equation*}
El sentido de la corriente debe hacer que $\vec L\times\vec B$ apunte hacia arriba. Si $\vec L=L\hat x$ y $\vec B=-B\hat z$, entonces
\begin{equation*}
    \vec L\times\vec B=LB\hat y,
\end{equation*}
que es justamente la direcci\'on necesaria para equilibrar el peso.

\newpage

\section*{Problema 10.}
Un conductor de masa $m=0.1\,\mathrm{kg}$ y largo $L=0.5\,\mathrm m$ descansa sin roce sobre un plano inclinado en $\alpha=10^\circ$. Hay un campo magn\'etico uniforme vertical de magnitud $B=5.2\,\mathrm{mT}$. Determinar la corriente necesaria para que el conductor permanezca en reposo.

\par\medskip\noindent Como no hay roce, el equilibrio importante es paralelo al plano. El peso tira hacia abajo del plano y la fuerza magn\'etica aporta una componente hacia arriba del plano.\par\medskip 


\begin{figure}[h!]
\centering
\begin{tikzpicture}[x=0.84cm,y=0.84cm,font=\small]
    \def\ang{22} % Angulo esquematico para separar visualmente las proyecciones.
    \draw[panel] (-1.28,-1.25) rectangle (16.10,8.78);

    % Plano inclinado y triangulo geometrico.
    \coordinate (A) at (0,0);
    \coordinate (B) at (13.65,0);
    \coordinate (C) at (13.65,{13.65*tan(\ang)});
    \draw[thick,black!80,fill=black!3] (A)--(B)--(C)--cycle;
    \draw[very thick,black!78] (A)--(C);
    \draw[thick,black!70] (13.00,0)--(13.00,0.65)--(13.65,0.65);

    % Rayado del plano.
    \foreach \t in {0.65,1.15,...,13.05}{
        \coordinate (H) at ($(A)!\t/13.65!(C)$);
        \draw[black!28,line width=0.35pt]
        ($(H)+({-0.09*sin(\ang)},{0.09*cos(\ang)})$)--($(H)+({0.18*cos(\ang)-0.09*sin(\ang)},{0.18*sin(\ang)+0.09*cos(\ang)})$);
    }

    % Angulo del plano con la horizontal.
    \draw[thick,black!70] (1.35,0) arc[start angle=0,end angle=\ang,radius=1.35];
    \node[etiqueta,fill=white,inner sep=1.5pt] at (1.98,0.31) {$\alpha$};

    % Conductor sobre el plano.
    \coordinate (S) at (6.95,{6.95*tan(\ang)});
    \coordinate (O) at ({6.95-0.46*sin(\ang)},{6.95*tan(\ang)+0.46*cos(\ang)});
    \begin{scope}[shift={(S)},rotate=\ang]
        \draw[thick,fill=white] (-0.82,0.00) rectangle (0.82,0.92);
        \draw[black!70,line width=1.20pt] (-0.66,0.46)--(0.66,0.46);
        \fill[black!75] (0,0.46) circle (0.07);
    \end{scope}

    % Marco local separado del DCL.
    \coordinate (E) at (2.12,{2.12*tan(\ang)+0.56});
    \draw[->,very thick,black]
        (E)--++({2.05*cos(\ang)},{2.05*sin(\ang)}) node[below right,fill=white,inner sep=1.5pt] {$\hat x$};
    \draw[->,very thick,black]
        (E)--++({-1.42*sin(\ang)},{1.42*cos(\ang)}) node[left,fill=white,inner sep=1.5pt] {$\hat y$};

    % Guias paralela y normal que pasan por el conductor.
    \draw[dashed,black!36,line width=0.75pt]
        ($(O)+({-2.65*cos(\ang)},{-2.65*sin(\ang)})$)--($(O)+({5.00*cos(\ang)},{5.00*sin(\ang)})$);
    \draw[dashed,black!36,line width=0.75pt]
        ($(O)+({1.85*sin(\ang)},{-1.85*cos(\ang)})$)--($(O)+({-2.30*sin(\ang)},{2.30*cos(\ang)})$);

    % Campo magnetico vertical.
    \draw[vecB,line width=1.20pt] (14.86,0.82)--(14.86,7.05);
    \node[orangevec,anchor=south] at (14.86,7.22) {$\vec B$ vertical};

    % Fuerza magnetica horizontal y proyeccion paralela.
    \draw[vecF,line width=1.15pt] (O)--++(3.45,0);
    \node[redvec,anchor=west,fill=white,inner sep=1.5pt] at ($(O)+(3.58,0.14)$) {$\vec F_B$};

    \draw[vecF,line width=1.05pt]
        (O)--++({3.05*cos(\ang)},{3.05*sin(\ang)});
    \node[redvec,anchor=south west,fill=white,inner sep=1.5pt] at ($(O)+({3.12*cos(\ang)+0.28},{3.12*sin(\ang)+0.38})$) {$F_{B,x}=F_B\cos\alpha$};
    \draw[dotted,redvec,line width=0.85pt]
        ($(O)+({3.05*cos(\ang)},{3.05*sin(\ang)})$)--($(O)+(3.45,0)$);
    \draw[redvec,thick] ($(O)+(1.08,0)$) arc[start angle=0,end angle=\ang,radius=1.08];
    \node[redvec] at ($(O)+(1.30,0.34)$) {$\alpha$};

    % Peso y componentes. Etiquetas fuera de los trazos principales.
    \draw[->,very thick,black!72] (O)--++(0,-3.28);
    \node[black!72,anchor=north,fill=white,inner sep=1.5pt] at ($(O)+(0,-3.38)$) {$\vec W=m\vec g$};

    \draw[->,thick,black!72]
        (O)--++({-2.52*cos(\ang)},{-2.52*sin(\ang)});
    \node[black!72,anchor=center,fill=white,inner sep=1.5pt] at (4.10,0.48) {$W_x=mg\sin\alpha$};

    \draw[->,thick,black!72]
        (O)--++({1.78*sin(\ang)},{-1.78*cos(\ang)});
    \node[black!72,anchor=west,fill=white,inner sep=1.5pt] at ($(O)+({1.78*sin(\ang)+0.22},{-1.78*cos(\ang)-0.02})$) {$W_y=mg\cos\alpha$};

    \draw[dotted,black!50,line width=0.72pt]
        ($(O)+({-2.52*cos(\ang)},{-2.52*sin(\ang)})$)--++({1.78*sin(\ang)},{-1.78*cos(\ang)});
    \draw[dotted,black!50,line width=0.72pt]
        ($(O)+({1.78*sin(\ang)},{-1.78*cos(\ang)})$)--++({-2.52*cos(\ang)},{-2.52*sin(\ang)});

    % Angulo entre el peso vertical y la direccion normal al plano.
    \draw[black!70,thick] ($(O)+(0,-0.82)$) arc[start angle=-90,end angle={-90+\ang},radius=0.82];
    \node[black!70] at ($(O)+(0.26,-1.00)$) {$\alpha$};

    % Normal.
    \draw[->,very thick,black]
        (O)--++({-1.72*sin(\ang)},{1.72*cos(\ang)});
    \node[black,anchor=east,fill=white,inner sep=1.5pt] at ($(O)+({-1.72*sin(\ang)-0.12},{1.72*cos(\ang)+0.06})$) {$\vec N$};

    % Nota breve, en una zona vacia del panel.
    \node[nota,anchor=north west,text width=9.40cm] at (0.40,8.30)
    {Se proyecta sobre $\hat x$.\\[-1pt]
    Geom\'etricamente: $F_{B,x}=F_B\cos\alpha$ y $W_x=mg\sin\alpha$.};
\end{tikzpicture}
\caption{Diagrama de cuerpo libre en el plano inclinado. El marco local se escoge paralelo y normal al plano para leer directamente las componentes.}
\end{figure}


\soluciontitulo

\par\medskip\noindent\textbf{Datos.}\par\smallskip
\begin{equation*}
    m=0.1\,\mathrm{kg},
    \qquad
    L=0.5\,\mathrm m,
    \qquad
    \alpha=10^\circ,
    \qquad
    B=5.2\,\mathrm{mT}=5.2\times10^{-3}\,\mathrm T.
\end{equation*}

\par\medskip\noindent\textbf{Modelo físico.}\par\smallskip
Elegimos un sistema local con $\hat x$ paralelo al plano inclinado y positivo hacia arriba del plano. La direcci\'on $\hat y$ es normal al plano.

La fuerza magn\'etica sobre el conductor tiene m\'odulo
\begin{equation*}
    |\vec F_B|=ILB.
\end{equation*}
En esta geometr\'ia, $\vec F_B$ queda horizontal. Como el plano forma un \~angulo $\alpha$ con la horizontal, la componente de $\vec F_B$ sobre el plano es
\begin{equation*}
    F_{B,x}=|\vec F_B|\cos\alpha.
\end{equation*}
El peso se descompone en
\begin{equation*}
    W_x=mg\sin\alpha,
    \qquad
    W_y=mg\cos\alpha.
\end{equation*}
La componente $W_x$ apunta hacia abajo del plano.

\par\medskip\noindent\textbf{Análisis.}\par\smallskip
Como no hay roce, basta equilibrar las fuerzas en la direcci\'on paralela al plano. Escogiendo positivo hacia arriba del plano,
\begin{equation*}
    F_{B,x}-W_x=0.
\end{equation*}
Por lo tanto,
\begin{equation*}
    |\vec F_B|\cos\alpha=mg\sin\alpha.
\end{equation*}
La direcci\'on normal solo determina el valor de $\vec N$, pero no cambia la corriente necesaria para el equilibrio paralelo al plano.

\par\medskip\noindent\textbf{Cálculos.}\par\smallskip
Sustituimos $|\vec F_B|=ILB$ y despejamos:
\begin{align*}
    ILB\cos\alpha&=mg\sin\alpha,\\
    I&=\frac{mg\sin\alpha}{LB\cos\alpha},\\
    I&=\frac{mg\tan\alpha}{LB}.
\end{align*}
Ahora se sustituyen los datos:
\begin{align*}
    I&=\frac{(0.1)(9.8)\tan(10^\circ)}{(0.5)(5.2\times10^{-3})}\,\mathrm A\\
     &=\frac{(0.98)(0.1763)}{2.6\times10^{-3}}\,\mathrm A\\
     &=\frac{0.1728}{2.6\times10^{-3}}\,\mathrm A\\
     &\approx66.46\,\mathrm A.
\end{align*}

\par\medskip\noindent\textbf{Resultado e interpretación física.}\par\smallskip
\begin{equation*}
    \boxed{I\approx6.6\times10^1\,\mathrm A}.
\end{equation*}
El sentido de la corriente debe escogerse para que la componente de $\vec F_B$ apunte hacia arriba del plano. Si se invierte la corriente, la fuerza magn\'etica apunta en el sentido contrario y el conductor no queda en equilibrio.

\par\medskip\noindent\textbf{Observación.}\enspace No se debe usar directamente $ILB=mg\sin\alpha$. La fuerza magn\'etica es horizontal, no paralela al plano; por eso aparece el factor $\cos\alpha$.\par\medskip 

\newpage

\section*{Problema 11.}
Una espira triangular rect\'angula tiene catetos $a=3\,\mathrm{cm}$ y $b=4\,\mathrm{cm}$. El plano de la espira forma un \~angulo $\alpha=45^\circ$ respecto de la direcci\'on $OY$. Circula una corriente $I=12\,\mathrm A$ y el campo magn\'etico uniforme es
\begin{equation*}
    \vec B=2.5\,\hat x\,\mathrm{mT}.
\end{equation*}
Calcular la fuerza magn\'etica sobre cada lado.

\par\medskip\noindent En cada tramo recto se usa $\vec F=I\vec\ell\times\vec B$. Lo delicado no es la f\'ormula, sino escribir bien cada vector $\vec\ell$ con el sentido de la corriente.\par\medskip 


\begin{figure}[h!]
\centering
\begin{tikzpicture}[
    x={(1.18cm,-0.32cm)},
    y={(0.72cm,0.48cm)},
    z={(0cm,1.18cm)},
    scale=1.07,
    font=\small
]
    \coordinate (O) at (0,0,0);
    \coordinate (A) at (0,0,3.20);
    \coordinate (C) at (1.78,2.85,0);

    % Ejes de referencia.
    \draw[->,thick,gray!68] (O)--(4.15,0,0) node[below] {$\hat x$};
    \draw[->,thick,gray!68] (O)--(0,3.92,0) node[right] {$\hat y$};
    \draw[->,thick,gray!68] (O)--(0,0,3.92) node[above] {$\hat z$};

    % Espira triangular.
    \draw[thick,black!84] (O)--(A)--(C)--cycle;
    \fill (O) circle (1.35pt) node[left] {$O$};
    \fill (A) circle (1.35pt) node[above left] {$A$};
    \fill (C) circle (1.35pt) node[right] {$C$};

    % Corriente en cada lado.
    \draw[vecV,line width=1.08pt] ($(A)!0.12!(O)$)--($(A)!0.40!(O)$) node[left,pos=0.55,fill=white,inner sep=1pt] {$I$};
    \draw[vecV,line width=1.08pt] ($(O)!0.15!(C)$)--($(O)!0.44!(C)$);
    \draw[vecV,line width=1.08pt] ($(C)!0.13!(A)$)--($(C)!0.42!(A)$);

    % Campo magnetico uniforme: flechas paralelas a x, fuera de la espira.
    \draw[vecB,line width=1.08pt] (-0.45,3.55,2.85)--(2.34,3.55,2.85);
    \node[orangevec,anchor=west,fill=white,inner sep=1.5pt] at (2.44,3.55,2.85) {$\vec B=B\hat x$};
    % Puntos medios para ubicar las fuerzas fuera de las aristas.
    \coordinate (Ma) at ($(O)!0.52!(A)$);
    \coordinate (Mb) at ($(O)!0.58!(C)$);
    \coordinate (Mh) at ($(C)!0.48!(A)$);

    % Fuerzas magneticas: perpendiculares a B y al tramo con corriente.
    \draw[vecF,line width=1.15pt] ($(Ma)+(0,-0.12,0)$)--++(0,-1.25,0);
    \node[redvec,anchor=east,fill=white,inner sep=1.5pt] at ($(Ma)+(0,-1.42,0.04)$) {$\vec F_a$};

    \draw[vecF,line width=1.15pt] ($(Mb)+(0.10,0.20,0.46)$)--++(0,0,-1.30);
    \node[redvec,anchor=north,fill=white,inner sep=1.5pt] at ($(Mb)+(0.10,0.20,-0.96)$) {$\vec F_b$};

    \draw[vecF,line width=1.15pt] ($(Mh)+(0.20,-0.35,0.10)$)--++(0,1.14,0.92);
    \node[redvec,anchor=south west,fill=white,inner sep=1.5pt] at ($(Mh)+(0.20,0.92,1.08)$) {$\vec F_h$};

\end{tikzpicture}
\caption{Espira triangular en un campo magn\'etico uniforme dirigido seg\'un $\hat x$. Las fuerzas se dibujan seg\'un la regla de la mano derecha.}
\end{figure}


\soluciontitulo

\par\medskip\noindent\textbf{Datos.}\par\smallskip
\begin{equation*}
    a=3\,\mathrm{cm}=0.03\,\mathrm m,
    \qquad
    b=4\,\mathrm{cm}=0.04\,\mathrm m,
    \qquad
    I=12\,\mathrm A,
\end{equation*}
\begin{equation*}
    \alpha=45^\circ,
    \qquad
    \vec B=2.5\times10^{-3}\hat x\,\mathrm T.
\end{equation*}

\par\medskip\noindent\textbf{Modelo físico.}\par\smallskip
Para cada lado rectil\'ineo:
\begin{equation*}
    \vec F=I\vec\ell\times\vec B.
\end{equation*}
Se toma el sentido de corriente indicado en la figura: baja por el lado vertical, avanza por el cateto $b$ y vuelve por la hipotenusa.

\par\medskip\noindent\textbf{Análisis geométrico.}\par\smallskip
Los vectores de longitud orientada son
\begin{align*}
    \vec\ell_a&=-a\hat z,\\
    \vec\ell_b&=b(\cos\alpha\,\hat y+\sin\alpha\,\hat x),\\
    \vec\ell_h&=a\hat z-b(\cos\alpha\,\hat y+\sin\alpha\,\hat x).
\end{align*}
La suma es nula, como debe ocurrir en una espira cerrada:
\begin{equation*}
    \vec\ell_a+\vec\ell_b+\vec\ell_h=\vec0.
\end{equation*}
Como $\vec B=B\hat x$, se usar\'an los productos
\begin{equation*}
    \hat z\times\hat x=\hat y,
    \qquad
    \hat y\times\hat x=-\hat z,
    \qquad
    \hat x\times\hat x=\vec0.
\end{equation*}

\par\medskip\noindent\textbf{Cálculos.}\par\smallskip
\textbf{Lado vertical.}
\begin{align*}
    \vec F_a
    &=I(-a\hat z)\times(B\hat x)\\
    &=-IaB(\hat z\times\hat x)\\
    &=-IaB\hat y.
\end{align*}

\textbf{Lado horizontal.}
\begin{align*}
    \vec F_b
    &=Ib(\cos\alpha\,\hat y+\sin\alpha\,\hat x)\times(B\hat x)\\
    &=IbB\cos\alpha(\hat y\times\hat x)+IbB\sin\alpha(\hat x\times\hat x)\\
    &=-IbB\cos\alpha\,\hat z.
\end{align*}

\textbf{Hipotenusa.}
\begin{align*}
    \vec F_h
    &=I\left[a\hat z-b(\cos\alpha\,\hat y+\sin\alpha\,\hat x)\right]\times(B\hat x)\\
    &=IaB\hat y+IbB\cos\alpha\,\hat z.
\end{align*}

Ahora se sustituyen los datos. Primero,
\begin{align*}
    IB&=(12)(2.5\times10^{-3})\\
      &=0.030\,\mathrm{N/m}.
\end{align*}
Luego,
\begin{align*}
    IaB&=(0.030)(0.03)=9.0\times10^{-4}\,\mathrm N,\\
    IbB\cos45^\circ&=(0.030)(0.04)\frac{\sqrt2}{2}
    \approx8.49\times10^{-4}\,\mathrm N.
\end{align*}

\par\medskip\noindent\textbf{Resultado e interpretación física.}\par\smallskip
\begin{align*}
    \boxed{\vec F_a=-9.0\times10^{-4}\hat y\,\mathrm N},\\
    \boxed{\vec F_b=-8.49\times10^{-4}\hat z\,\mathrm N},\\
    \boxed{\vec F_h=9.0\times10^{-4}\hat y\,\mathrm N+8.49\times10^{-4}\hat z\,\mathrm N}.
\end{align*}
Como chequeo,
\begin{equation*}
    \vec F_a+\vec F_b+\vec F_h=\vec0.
\end{equation*}
Esto es coherente con una espira cerrada dentro de un campo magn\'etico uniforme: la fuerza total se anula, aunque puede quedar torque sobre la espira.

\end{document}
